%
%============= CHAPTER BREAK =================%
%
\chapter{Vettius Valens (152--165 CE)}
\vspace*{3mm}%
\justifying%
%
\noindent{}Born and raised in Antioch, %
Vettius Valens learned and wrote about astrological traditions in Alexandria, Egypt. %
Next to Ptolemy, Valens was the best known and most highly regarded astrologer of the %
2nd century CE. %
According to Neugebauer (1954), Valens wrote his astrological compendium, %
the \textit{Anthology}, largely between 152 and 165 CE. \\[1.5mm] %%

\noindent{}\textbf{ἀφελότης} / \textit{aphelotēs} / “simplicity” first occurs in non-Christian literature %
in Valensʼ \textit{Anthology}, in two different passages. %

\vspace*{3mm} 
 \begin{quote}\textls[-15]{ἣν εἴ τις φυλάττοι, 
 εὐκατάληπτον ἕξει τὴν τῶν ζητουμένων ἀλήθειαν. %
 τινὰ μὲν οὖν καὶ μυστικῶς διεγράψαμεν εἰς διάκρισιν %
 καὶ σύγκρισιν τῶν ἐντυγχανόντων οὐ φθόνῳ φερόμενοι οὐδὲ \textbf{ἀφελότητι}, 
 ἀλλὰ προθυμίαν τινὰ καὶ ἐπιμονὴν βουλευόμενοι παρεισφέρειν τῷ μαθήματι· 
 πᾶς γὰρ καταθυμίου προαιρέσεως ἀκολάστως τυχὼν ἀχάριστον τὴν δωρεὰν ἡγήσατο, 
 ὁ δὲ μετὰ πόνου καὶ ζητήσεως οὐ μόνον ἡδονὴν %
 ἀλλὰ καὶ κατόρθωσιν τῇ πραγματείᾳ ἐπεισήνεγκεν.} (\textit{Anth}. 3.10) \\[1mm] %
 If anyone should guard this [method], %
 from things sought he will hold the truth easily recognizable. %
 While then some things we delineated mystically for differentiation and comparison %
 of things encountered, not being carried away by envy or by \textbf{simplicity}, %
 but being desirous to introduce a certain readiness and endurance to the study. %
 For everyone who chances effortlessly on an heartfelt preference considers the gift ungrateful, %
 but one who with toil and seeking brings not only pleasure, %
 but also accomplishment as well in the research. \end{quote} % 
 \vspace*{1mm}

 \begin{quote}πολλάκις γὰρ εὑρίσκεταί τις ὢν δημιουργὸς καὶ μεθοδικὸς %
 πρὸς τὰς ἐπιβολὰς τῶν πράξεων καὶ ἐπιτευκτικὸς κἂν ἄπειρος γραμμάτων ἐπιτυγχάνῃ, %
 ὁ δὲ πάνυ πεπαιδευμένος ῥᾳδίως ἁλίσκεται ὡς ἄπειρος τῶν πρακτέων, %
 ὑπ’ \textbf{ἀφελότητος} καὶ ἀδιοικησίας προδεδομένος· %
 καὶ οὗτος ὀδυνώμενος ματαίαν ἡγεῖται τὴν τῆς παιδείας %
 ἐπιβολὴν καὶ εὐδαίμονα προκρίνει τὸν ἀμαθῆ. (\textit{Anth}. 6.1) \\[1mm] %
 For oftentimes someone may be found who, %
 being industrious and methodical in their endeavors of action and effective, %
 even if he chances to be unacquainted with writings, %
 while one who is altogether educated easily succumbs as one unacquainted with what must be done, %
 betrayed by \textbf{simplicity} and want of management. %
 And this one who is grieved considers empty the endeavor of education %
 and prefers the unlearned as fortunate.\end{quote} % 
 \vspace*{3mm}

\noindent{}In both passages, Valens uses the lemma in a negative sense. %
He pictures “simplicity” as an external and seductive force in the life of the learned, %
akin to the lure of laziness and naïveté. %
And yet in the same context he also recognizes the potential virtue of the unlearned, %
even though he does not directly associate the word “simplicity” with that virtue. %

These two examples from Valens are quite reasonable as background for %
the usage in Acts, %
where the term ἀφελότης / \textit{aphelotēs} / “simplicity” is an NT \textit{hapax legemenon}. %
Here it seems that Acts presumes yet also inverts the sense, %
turning “simplicity of heart” into a virtue. %
This shift may owe its inspiration to a secondary source, %
canonical Ephesians, which in a verse apparently missing from Marcionite Laodiceans, %
addresses slaves and encourages them to act with “singleness of heart” / % 
ἁπλότητι τῆς καρδίας (Eph 6:5). 

\vspace*{3mm} 
 \begin{quote}καθ’ ἡμέραν τε προσκαρτεροῦντες ὁμοθυμαδὸν ἐν τῷ ἱερῷ,
 κλῶντές τε κατ’ οἶκον ἄρτον, %
 μετελάμβανον τροφῆς ἐν ἀγαλλιάσει καὶ \textbf{ἀφελότητι} καρδίας, %
 ⁴⁷ αἰνοῦντες τὸν θεὸν καὶ ἔχοντες χάριν πρὸς ὅλον τὸν λαόν. (Acts 2:46--47) \\[1mm] %
 Daily with one accord in the temple, %
 breaking bread in each home, %
 partaking of food with delight and \textbf{simplicity} of heart, %
 ⁴⁷ praising the God and having favor with the entire people.\end{quote} % 
 \vspace*{3mm} 

In the subsequent context, %
Peter gives a formal, public speech to an audience that includes %
“the priests and the military commander of the temple and the Sadducees” / %
οἱ ἱερεῖς καὶ ὁ στρατηγὸς τοῦ ἱεροῦ καὶ οἱ Σαδδουκαῖοι (Acts 4:1). %
After hearing Peter speak, %
“they were amazed” / ἐθαύμαζον, “because they were unlettered and ignorant men” / %
ὅτι ἄνθρωποι ἀγράμματοί εἰσιν καὶ ἰδιῶται (4:11). %
Acts thus evokes and challenges the very contrast that Valens described %
between the learned and unlearned, %
perhaps subtly depicting Peter as a second Aesop, %
suddenly and divinely gifted with oratorical skill and boldness. %

The lemma also appears in manuscripts of the late 2nd century sermon \textit{On the Pasch} by Melito of Sardis. %
These concluding lines could be read as Melitoʼs closing, poetic benediction, %
or as an editorial \textit{explicit}.\endnote{%
A summary conclusion often used in manuscripts to distinguish %
the previous work from the following one.} %

\vspace*{3mm} 
 \begin{quote}⸤Εἰρήνη τῷ γράψαντι καὶ τῷ ἀναγινώσκοντι καὶ τοῖς ἀ⸥- \newline
 ⸤γαπῶσι τὸν κύριον ἐν \textbf{ἀφελότητι} καρδίας⸥. (\textit{Pascha}, 825--826) \\[1mm] %
 Peace to him who writes and him who reads and them who \newline 
 love the Lord in \textbf{simplicity} of heart. \end{quote} % 
 \vspace*{3mm}

In the TLG, for both Christian and non-Christian literature, %
the lemma next occurs in the 4th century CE. %
It also seems to be absent from Greek papyri and inscriptions, %
strengthening the link of this rare word between Valens and canonical Acts. \clearpage %

\noindent{}\textbf{εὐμετάδοτος} / \textit{eumetadotos} / “readily giving” %
first appears in a chronologically secure text %
in the \textit{Anthology} of Valens. %
Most of the eleven (!) occurrences of this lemma are found %
within two clusters in book one %
(1.2.50, 1.2.164, 1.2.197, 1.20.42, 1.20.161, 1.20.167, 1.20.230), %
a volume written early in that period. %
Only one hit appears in book 2 (2.21.7), %
but two more in the appendix for book 1 (10.1.149, 10.1.172). %

Most of the hits, as might be expected, pertain to the significance of various stars %
and constellations as keys to interpret the meaning of celestial events, %
the characteristics of human souls and times of birth, %
and the dispositions of those born under various astrological signs. %
All of these employ the term as a positive character attribute. %

Marcus Aurelius, in his Stoic-adjacent \textit{Meditations}, %
subsequently uses this lemma twice, both times as a positive character attribute, %
one apparently associated more with philosophy than astrology (though Stoics admittedly often held these together): % 

\vspace*{2mm} 
 \begin{quote}καὶ ἔτι παρ’ αὐτοῦ τὸ ὁμαλὲς καὶ εὔτονον ἐν τῇ τιμῇ τῆς φιλοσοφίας· %
 καὶ τὸ εὐποιητικὸν καὶ τὸ \textbf{εὐμετάδοτον} ἐκτενῶς καὶ τὸ εὔελπι· \\[1mm] %
 And also from him [i.e., Severus, I learned] what is equitable %
 and energizing in the value of philosophy, %
 and what is beneficent, and what is \textbf{readily giving} extensively, and what is hopeful. (\textit{Medit}. 1.14.2) %
 \end{quote} % 
\vspace*{1mm}
 
 \begin{quote}Ὅταν εὐφρᾶναι σεαυτὸν θέλῃς, ἐνθυμοῦ τὰ προτερήματα τῶν
 συμβιούντων· οἷον τοῦ μὲν τὸ δραστήριον, τοῦ δὲ τὸ αἰδῆμον,
 τοῦ δὲ τὸ \textbf{εὐμετάδοτον}, ἄλλου δὲ ἄλλο τι. οὐδὲν γὰρ οὕτως
 εὐφραίνει ὡς τὰ ὁμοιώματα τῶν ἀρετῶν ἐμφαινόμενα τοῖς ἤθεσι
 τῶν συζώντων καὶ ἀθρόα ὡς οἷόν τε συμπίπτοντα. (\textit{Medit}. 6.48.1) \\[1mm] %
 Whenever you wish to cheer yourself up, %
 take to heart the superior qualities of those living with you, %
 such as the efficacy of this one, or the modesty of that one, %
 or the \textbf{ready generosity} of that one, or something else of another. %
 For nothing cheers up as much as the likenesses of the virtues %
 shining in the ethics of those living with you and %
 shared together as such and co-occuring. 
  \end{quote} % 
 \vspace*{2mm}

 In its single appearance in the Pastorals, the term carries a decidedly ethical sense. %
 This makes the emperorʼs philosophical reflections the more likely, immediate conduit %
 rather than the astrology of Valens, even if Valens coined or first popularized the term. %
 The focus on the practice of virtue and modesty in community also makes %
 Pauline proto-Orthodox guidance resonate with imperial introspection (1Tim 6:17--18): %

\vspace*{2mm} 
 \begin{quote}Τοῖς πλουσίοις ἐν τῷ νῦν αἰῶνι παράγγελλε μὴ %
 ὑψηλοφρονεῖν μηδὲ ἠλπικέναι ἐπὶ πλούτου ἀδηλότητι, %
 ἀλλ’ ἐπὶ θεῷ τῷ παρέχοντι ἡμῖν πάντα πλουσίως εἰς ἀπόλαυσιν, % 
 ¹⁸ ἀγαθοεργεῖν, πλουτεῖν ἐν ἔργοις καλοῖς, \textbf{εὐμεταδότους} εἶναι, κοινωνικούς. \\[1mm] %
 To those rich in this present age, command them not to %
 be high-minded nor to hope in uncertain wealth, %
 but instead [to hope] in God who supplies us everything richly for enjoyment, % 
 ¹⁸ [and] to do good, to be rich in beautiful deeds, %
 to be \textbf{readily giving}, sharing. %
 \end{quote} %
 \vspace*{2mm} 

\noindent{}The character trait of ready generosity, %
initially indicated by Vettiusʼ stars, %
later transformed by the Stoic philosopher-emperor into philosophical aspiration, %
now becomes a hallmark of early-Orthodox ethical guidance in the form of \textit{synkrisis}: %
litanies of virtues and vices composed in the interest of character and community formation. \\[1.5mm] %

\noindent{}\textbf{ἰσάγγελος} / \textit{isangelos} / “angel-like” %
has 553 TLG hits, yet it seems to have first become popular in the mid-2nd century. %
It first appears in a non-Christian writing in an appendix to Valensʼ \textit{Anthology}. %
The section begins with the heading, %
Περὶ σχηματισμοῦ τῶν ἀστέρων κατὰ σύνοδον γινομένου οὐχ ἑνὸς πρὸς
ἕνα ὡς διείληπται πρότερον ἀλλ’ ἑνὸς πρὸς δύο ἢ καὶ πλείους / %
“On the Configuration of the Stars according to Meeting, %
Not When One Comes unto One, as Was Explained Before, %
but When One Comes unto Two or even Many”. %

\vspace*{2mm} 
 \begin{quote}Ὁ δὲ τοῦ Ἡλίου καὶ Ἑρμοῦ καὶ Σελήνης σεμνούς, καθαρίους, εὐυποκρίτους, %
 οἰκονομικούς, μυστηρίων μετόχους, κατορθωτὰς πραγμάτων,
 πλείστην φαντασίαν τῆς ὑπάρξεως κεκτημένους, σωματοφύλακας,
 \textbf{ἰσαγγέλους}, ἐπάνω χρημάτων, γραμμάτων, ψήφων τεταγμένους· %
 τῶν δὲ τοιούτων καὶ ὁ λόγος ἐξισχύσει πρὸς συμβουλίαν καὶ ἡ διδαχή. (\textit{Anth}. App. 10.1, lines 232--236)\endnote{%
 Pingree 1986:413.} \\[1mm] %
\textls[-15]{ The co-configuration of the Sun and Mercury and the Moon [renders persons who are] %
 reverent, pure, part-suited, house-managerial, partakers of mysteries, %
 accomplished in dealings, having been possessed of the greatest prestige of substance, %
 bodyguards, \textbf{angel-equals}, being stationed over means, letters, and votes. %
 The word and the teaching of such ones prevail on the assembly.} \end{quote} % 
 \vspace*{2mm}

This same word appears in a close parallel shared between Marcionʼs \textit{Evangelion} and canonical Luke, %
with subtle surrounding differences between the two. %
Here we only quote our Marcionite reconstruction:\endnote{Cp. SGC at Lk2 20:35--36. %
Tertullian attests: “Those whom god has made worthy of that age, %
for inheritance and resurrection from the dead, %
neither marry nor are married, %
for neither yet do they die, since they are equals of angels of God, %
having been made sons of resurrection” / %
\textit{quos vero dignatus sit deus illius aevi possessione et resurrectione a mortuis neque nubere neque nubi, %
quia nec morituri iam sint, cum similes angelorum sint dei, %
resurrectionis filii facti} (\textit{Marc}. 4.38.5); %
“Whom yet god has made worthy of that age, to God they join themselves” / %
\textit{quos autem dignatus est deus illius aevi, deo adiungant} (\textit{Marc}. 4.38.7.)} %

\vspace*{2mm} 
 \begin{quote}οὓς κατηξίωσεν ὁ θεὸς τοῦ αἰῶνος ἐκείνου τῆς κληρονομίας %
 καὶ τῆς ἀναστάσεως τῆς ἐκ νεκρῶν οὔτε γαμοῦσιν οὔτε γαμίσκονται\endnote{%
 See chapter 2, infra, s.v., γαμίζω, on this lemma choice.} %
 ³⁶ οὐδὲ γὰρ ἀποθανεῖν, \textbf{ἰσάγγελοι} γάρ εἰσιν καὶ %
 υἱοί εἰσιν θεοῦ τῆς ἀναστάσεως υἱοὶ ὄντες. \\[0.25mm] %
 Those whom the God made worthy of the age, %
 that of the inheritance and of the resurrection that [is] from the dead, %
 neither do they marry, nor are they married, %
 ³⁶ nor do they die, for they are \textbf{angel-equals} and %
 are sons of God, being sons of the resurrection. %
 \end{quote} % 
\clearpage

The Gospel passages share with Valens not only the new lemma “angel-equal”, %
but also a focus on divine destiny, possessions, authority, and royal responsibilities. %
By my lights, Valens is a likely source behind both Marcionʼs \textit{Evangelion} and canonical Luke at this point. %
It is usually assumed that Justin quotes canonical Luke 20:25, %
but the quotation shares elements with Marcionʼs \textit{Evangelion} and canonical Luke, %
while also evidencing distinctive verb and noun choices: %

\vspace*{2mm} 
 \begin{quote}ὅπερ καὶ ὁ κύριος ἡμῶν εἶπεν, %
 ὅτι Οὔτε γαμήσουσιν οὔτε γαμηθήσονται, %
 ἀλλὰ \textbf{ἰσάγγελοι} ἔσονται, τέκνα τοῦ θεοῦ τῆς ἀναστάσεως ὄντες. (\textit{Tryph}. 81.4) \\[1mm] %
 Those whom the Lord of ours said, %
 “Neither will they marry nor will they be made to marry, %
 but they will be \textbf{angel-equals}, being children of the God of the resurrection.”\end{quote} % 
 \vspace*{2mm} %

\noindent{}In other early Christian literature, the lemma %
also appears in Ps-Clem, \textit{Hom}. 17.16.5, %
which speaks of bodies becoming light and “\textbf{angel-equals}” / \textbf{ἰσάγγελοι}, %
akin to the way angels of light become flesh to be seen by humans. %
It also appears numerous times in the writings of ClemAlex.\endnote{%
E.g., \textit{Paed}. 1.6.36.5; \textit{Strom}. 6.13.105.1, 7.10.57.5, etc.}

Justinʼs \textit{Dialogue with Trypho} could date anywhere from 155 to 167 CE.\endnote{%
Eusebius claimed Justin wrote this Socratic dialogue in Ephesus shortly after Hadrianʼs reign, %
but that claim is viewed as completely unreliable in recent scholarship. %
See Bagatti 1979, Hamman 1995, and Slosser 2003:xii. Recent scholars concur %
that \textit{Dial.Trypho} %
refers back to his \textit{I Apol}. %
But there is some quibbling and ambiguity over the exact year(s) of both texts: %
Minns \& Parvis (2009:33) 153 CE and sometime subsequent; %
Allert (2005:32--34) 151--154 CE and 155--167 CE; %
Rokéah (2002:2) ca 155 CE and ca 160 CE.} %
The lemma is not apparently attested in any Greek papyri or Greek inscriptions %
digitized in the major databases. %
By my lights, if a progenitor is to be chosen from among available candidates, %
then Vettius Valens is the strongest. %
On this reading, Marcionʼs \textit{Evangelion} (at least in its secondary redaction), %
Justin, and a redactor of canonical Luke obtained this ostensibly dominical teaching %
not from Jesus of Nazareth, but from their famous astrologer colleague in Rome. %
\clearpage

\noindent{}\textbf{λογομαχία} / \textit{logomachia} / “word-fight” %
has its first chronologically secure TLG hit in a fragment by Valens about the “natal star” / γενεθλίου.\endnote{%
Cumont 1929:167.} %

\vspace*{3mm} 
 \begin{quote}ἐὰν δὲ ἔλθῃ εἰς τὸ τετράγωνον τῆς γενέσεως, %
 ποιεῖ \textbf{λογομαχίας} πρὸς θηλυκὰ πρόσωπα καὶ %
 βιωτικὰς βλάβας καὶ σωματικὰς ἐπιζητήσεις. \\[1mm] % 
 \textls[-15]{Now if it comes to the quadrangle of the birth, 
 it indicates \textbf{word-fights} after a womanish character and %
 biological harms and bodily cravings.}\end{quote} % 
\vspace*{3mm}

\noindent{}Its single NT appearance is in the Pastorals, specifically 1Tim 6:3--5: 

\vspace*{3mm} 
 \begin{quote}\textls[-16]{εἴ τις ἑτεροδιδασκαλεῖ καὶ μὴ προσέρχεται ὑγιαίνουσιν λόγοις, %
 τοῖς τοῦ κυρίου ἡμῶν Ἰησοῦ Χριστοῦ, καὶ τῇ κατ’ εὐσέβειαν διδασκαλίᾳ, %
 ⁴ τετύφωται, μηδὲν ἐπιστάμενος, %
 ἀλλὰ νοσῶν περὶ ζητήσεις καὶ \textbf{λογομαχίας}, %
 ἐξ ὧν γίνεται φθόνος, ἔρις, βλασφημίαι, ὑπόνοιαι πονηραί, %
 ⁵ διαπαρατριβαὶ διεφθαρμένων ἀνθρώπων τὸν νοῦν καὶ ἀπεστ\-ερημένων τῆς ἀληθείας, %
 νομιζόντων πορισμὸν εἶναι τὴν εὐσέβειαν.} \\[1mm] %
 If anyone teaches otherwise and does not proceed to speak with healthy words, %
 with those of our Lord Jesus Christ, and in teaching according to devotion, %
 ⁴ he has become demented, understanding nothing, %
 but instead diseased over cravings and \textbf{word-fights}, %
 from which comes malice, envy, blasphemies, evil suspicions, %
 ⁵ constant wrangling between humans corrupted in the mind %
 and defrauding of the truth, who deem monetary gain to be devotion. %
\end{quote} % 
 \vspace*{3mm}

\noindent{}The echo of the new and rare word λογομαχία / \textit{logomachia} is notable, %
especially given that here it is also plural and accusative. %
Supplemental confirmation of dependence obtains in the adjacent use of lemmas %
for “craving”, differentiated by a mere prefix. Compare ἐπιζητήσεις / \textit{epidzētēseis} %
and ζητήσεις / \textit{dzētēseis}. %
The semantics of health (good and/or ill) also overlap. %
It appears that even when the Stoic emperor did not explicitly appropriate the terminology %
of Valens --- at least based on the extant Greek texts of Aurelius available in the TLG ---, %
another of his distinctive neologisms still found its way %
into early-Orthodox ethical instruction, this time as a sign of shame. %